\section{Ответы на теоретические контрольные вопросы}

\begin{enumerate}
    \item
    \textit{Дайте определение выпуклого множества. Сформулируйте
    определение своими словами и приведите 2--3 примера выпуклых и
    невыпуклых множеств на плоскости.}

    Определение~\ref{dfn:convex_set} можно понимать так: выпуклое множество содержит отрезок между любой парой своих точек.
    Выпуклыми примерами служат круг, квадрат и полуплоскость.
    Невыпуклыми являются кольцо, звездная область и многоугольник
    с выемкой.

    \item
    \textit{Сформулируйте и докажите свойство сохранения выпуклости при
    пересечении множеств. Почему это свойство важно для практических
    приложений?}

    Согласно теореме~\ref{thm:intersection}, пересечение выпуклых
    множеств выпукло, потому что для любой пары точек пересечения
    отрезок между ними одновременно лежит в каждом из исходных множеств.
    На практике это важно при формировании области допустимых решений
    системы ограничений.

    \item
    \textit{Объясните геометрический смысл выпуклой линейной комбинации
    точек. Как связаны понятия выпуклой оболочки и выпуклых комбинаций?}

    Коэффициенты выпуклой комбинации являются неотрицательными весами,
    сумма которых равна единице. Поэтому результат является взвешенным
    средним исходных точек и находится в их выпуклой оболочке. Все точки
    выпуклой оболочки можно получить таким способом.

    \item
    \textit{Какие методы проверки выпуклости многоугольника вы знаете?
    Опишите алгоритм на основе векторных произведений.}

    Для проверки выпуклости можно использовать проверку углов и знаков
    векторных произведений, ориентационный тест, проверку пересечений
    ребер, а также построение выпуклой оболочки с последующим сравнением.

    В работе применен линейный проход по последовательным тройкам вершин
    с фиксацией первого ненулевого знака векторного произведения. Если
    среди последующих произведений встречается противоположный знак,
    многоугольник является невыпуклым.

    \item
    \textit{В чем заключается практическая значимость выпуклых множеств
    в машинном обучении и оптимизации?}

    Выпуклость упрощает оптимизацию: у выпуклой задачи локальный минимум
    является глобальным, а область допустимых решений не содержит
    раздельных локальных подобластей.

    Выпуклые множества используются в линейном и квадратичном
    программировании, методах опорных векторов, регуляризации и анализе
    ограничений машинного обучения.
\end{enumerate}
